\newpage
\phantomsection
\label{prf:metric-diameter-monotone-under-inclusion}
\LRAProofFor{thm:metric-diameter-monotone-under-inclusion}

\begin{remark*}[Return]
\hyperref[thm:metric-diameter-monotone-under-inclusion]{Return to Theorem}
\end{remark*}

\begin{theorem*}[Diameter is Monotone Under Inclusion]
Let \((X,d)\) be a metric space and let \(A,B\subseteq X\).  If
\(A\subseteq B\), \(D(A)\) is nonempty, and \(D(B)\) is bounded above, then
\(\operatorname{diam}(A)\leq\operatorname{diam}(B)\).
\end{theorem*}

\begin{proof}[Professional Standard Proof]
\LRAProofBodyStart
By \hyperref[thm:metric-diameter-set-monotone]{monotonicity of diameter
distance sets}, \(D(A)\subseteq D(B)\).  Since \(D(B)\) is bounded above,
every upper bound for \(D(B)\) is also an upper bound for \(D(A)\).  The least
upper bound of \(D(A)\) is therefore no larger than the least upper bound of
\(D(B)\).  Hence
\[
  \operatorname{diam}(A)=\sup D(A)\leq \sup D(B)=\operatorname{diam}(B).
\]
\end{proof}

\begin{proof}[Detailed Learning Proof]
\LRAProofBodyStart
The diameter is defined as a supremum:
\[
  \operatorname{diam}(A)=\sup D(A),
  \qquad
  \operatorname{diam}(B)=\sup D(B).
\]
The inclusion \(A\subseteq B\) implies
\[
  D(A)\subseteq D(B),
\]
because every distance between two points of \(A\) is also a distance between
two points of \(B\).

Now use the order property of suprema.  Since \(D(B)\) is bounded above and
\(D(A)\) is nonempty, the supremum \(\sup D(A)\) is controlled by any upper
bound for \(D(B)\).  In particular, \(\sup D(B)\) is an upper bound for
\(D(B)\), hence also an upper bound for \(D(A)\).  Because \(\sup D(A)\) is
the least upper bound of \(D(A)\), it must satisfy
\[
  \sup D(A)\leq \sup D(B).
\]
Translating back to diameters gives
\[
  \operatorname{diam}(A)\leq \operatorname{diam}(B).
\]
\end{proof}

\begin{remark*}[Proof structure]
First reduce the claim to \(D(A)\subseteq D(B)\).  Then apply the elementary
supremum principle that the supremum of a nonempty subset is at most the
supremum of a bounded-above containing set.
\end{remark*}

\begin{dependencies}
\begin{itemize}
  \item \hyperref[def:metric-diameter]{Definition: Diameter}
  \item \hyperref[thm:metric-diameter-set-monotone]{Theorem: Diameter Distance Sets are Monotone}
\end{itemize}
\end{dependencies}

\clearpage
